%%%%%%%%%%%%%%%%%%%% author.tex %%%%%%%%%%%%%%%%%%%%%%%%%%%%%%%%%
%
% sample root file for your "contribution" to a contributed volume
%
% Use this file as a template for your own input.
%
%%%%%%%%%%%%%%%% Springer Nature %%%%%%%%%%%%%%%%%%%%%%%%%%%%%%%%


% RECOMMENDED %%%%%%%%%%%%%%%%%%%%%%%%%%%%%%%\documentclass[runningheads]{llncs}
\documentclass[conference,compsoc]{IEEEtran}

\raggedbottom

\ifCLASSOPTIONcompsoc
  % IEEE Computer Society needs nocompress option
  % requires cite.sty v4.0 or later (November 2003)
  \usepackage[nocompress]{cite}
\else
  % normal IEEE
  \usepackage{cite}
\fi


\usepackage[T1]{fontenc}
\usepackage{graphicx}
\usepackage{booktabs}
\usepackage{microtype}
\usepackage{pgfplots}
\usepackage{pgfplotstable}
\usepackage{float}
\usepackage{caption}
\usepackage{titlesec}
\usepackage{multirow}
\usepackage{url}
% \setcounter{secnumdepth}{3}
% \titlespacing*{\section}
% {3pt}{1.2ex plus .2ex}{0.8ex}

% \titlespacing*{\subsection}
% {3pt}{1ex}{0.5ex}

% \usepgfplotslibrary{statistics}
% \pgfplotsset{compat=1.18}
\usepackage{pgfplots}
\usepgfplotslibrary{statistics}
\pgfplotsset{compat=1.18}

\def\BibTeX{{\rm B\kern-.05em{\sc i\kern-.025em b}\kern-.08em
    T\kern-.1667em\lower.7ex\hbox{E}\kern-.125emX}}

\makeatletter
\newcommand{\linebreakand}{%
  \end{@IEEEauthorhalign}
  \hfill\mbox{}\par
  \mbox{}\hfill\begin{@IEEEauthorhalign}
}
\makeatother

\hyphenation{op-tical net-works semi-conduc-tor}

\begin{document}
% \motto{This is my motto.}
\title{HiCWE: Leveraging Historical CWE Exploitability and Bi-Encoder-based CWE Classification for Early Vulnerability Prioritization}

\author{

\IEEEauthorblockN{Person A}
\IEEEauthorblockA{
    \textit{Some University} \\
    email@email.com
}

\and
\IEEEauthorblockN{Person B}
\IEEEauthorblockA{
    \textit{Some University} \\
    email@email.com
}

\and
\IEEEauthorblockN{Person C}
\IEEEauthorblockA{
    \textit{Some University} \\
    email@email.com
}

\and
\IEEEauthorblockN{Person D}
\IEEEauthorblockA{
    \textit{Some University} \\
    email@email.com
}
}
% \author{
% Kobra Khanmohammadi\inst{1}\orcidID{0009-0004-1414-2111}  \and
% Alexander Da Re\inst{1} \and Benjamin Whelan\inst{1} \and
% Raphaël Khoury\inst{2}\orcidID{0000-0002-7625-3384} 
% }

% \authorrunning{K. Khanmohammadi et al.}

% \institute{
% Sheridan College, Ontario, Canada\\
% \email{kobra.khanmohammadi@sheridancollege.ca, dareal@sheridancollege.ca, whelaben@sheridancollege.ca}
% \and
% Université du Québec en Outaouais (UQO), Canada\\
% \email{raphael.khoury@uqo.ca}}
%
% Use the package "url.sty" to avoid
% problems with special character[s
% used in your e-mail or web address
%
\maketitle


%%%%%%%%%%%%%%%%%%%%%%%
%%%%
\begin{abstract}
Given limited security resources and the growing
number of disclosed vulnerabilities, organizations
must prioritize vulnerabilities for remediation.
Although the National Vulnerability Database (NVD)
provides severity assessments, vulnerability severity
does not necessarily reflect the likelihood of
exploitation in the wild. In this paper, we investigate
the relationship between real-world exploitation
and Common Weakness Enumeration (CWE) categories,
which characterize the underlying root causes of
software vulnerabilities. We propose a framework
that prioritizes newly disclosed vulnerabilities
based on the historical exploitation rates of CVEs
associated with their CWE categories. We analyze
192,788 CVEs across 768 CWEs, including 3,097
exploited vulnerabilities, to assess the temporal
stability of CWE-level exploitation patterns.
To address missing or delayed CWE assignments,
we introduce HiCWE, a hierarchy-aware bi-encoder
that automatically maps CVE descriptions to CWE
definitions, achieving 75.3\% top-1 classification
accuracy. We further evaluate historical CWE-based
prioritization using official NVD CWE assignments
and rolling temporal evaluations from 2021 to 2025,
achieving ROC-AUC scores between 0.692 and 0.727.
Our findings establish historical CWE exploitability
as a useful complementary signal and provide a
framework for incorporating root-cause information
into early vulnerability prioritization.

\end{abstract}

\begin{IEEEkeywords}
Common Vulnerabilities and Exposures (CVE),
Common Weakness Enumeration (CWE),
Large Language Models,
Vulnerability Prioritization,
Exploitability Analysis
\end{IEEEkeywords}
%%%%%%%%%%%%%%%%%%%%%%%%%%%%%%%%%%%%%%%%%%%%

\IEEEpeerreviewmaketitle


\section{Introduction}
\label{sec:introduction}

The increasing number of publicly disclosed software
vulnerabilities makes effective vulnerability
prioritization essential. Security teams must identify
which vulnerabilities require immediate remediation,
often before sufficient information about their
exploitation becomes available. Although practitioners
widely use the Common Vulnerability Scoring System
(CVSS) to assess severity, previous studies have shown
that severity alone does not reliably predict
real-world exploitation
\cite{jacobs2023epss,sherif2026kri,jiang2025survey}.
To address this limitation, Jacobs et
al.~\cite{jacobs2023epss} developed the Exploit
Prediction Scoring System (EPSS), which estimates
exploitation likelihood using vulnerability
characteristics and threat-related information.
However, some of these indicators may be unavailable
or incomplete immediately after disclosure,
limiting their usefulness for early prioritization.

One potential complementary signal is the Common
Weakness Enumeration (CWE), a taxonomy maintained
by MITRE that categorizes software weaknesses
according to the underlying root causes of
vulnerabilities. Unlike CVEs, which identify
individual vulnerabilities, CWEs represent recurring
weakness types across software products.
Previous studies have examined relationships
between software weaknesses and exploitation\cite{mell2020suite, galhardo2020measurements} 
while MITRE's CWE Top 25 identifies particularly
dangerous weakness categories~\cite{cwe_top25}.
However, these analyses do not establish whether
historical CWE exploitation rates remain sufficiently
stable over time to predict the exploitation of
future vulnerabilities. We investigate this gap
by examining the temporal consistency and
predictive usefulness of historical weakness-level
exploitation patterns.

A second challenge is the availability of CWE
assignments. Newly disclosed CVEs may initially
lack specific CWE mappings, preventing security
teams from immediately using historical
weakness-level information. Researchers have
proposed automated CVE-to-CWE classification
methods using BERT~\cite{das2021v2w},
TF-IDF-based similarity~\cite{albanese2024cve2cwe},
and transformer-based or generative language
models~\cite{mosievskiy2026roberta,
simonetto2026knowing}. Nevertheless, short
and ambiguous CVE descriptions, overlapping
CWE definitions, and hierarchical relationships
between weaknesses continue to complicate
classification. Moreover, fixed-label classifiers
cannot naturally accommodate new CWE definitions
without modifying their prediction layers.
These limitations motivate a hierarchy-aware
semantic matching approach.

To address these challenges, we propose a
three-phase framework for early vulnerability
prioritization. First, we analyze historical
CVE--CWE mappings and exploitation evidence
to assess the stability and predictive value
of CWE-level exploitation patterns.
Second, we introduce HiCWE, a hierarchy-aware
bi-encoder that uses contrastive learning
and CWE taxonomy relationships to map CVE
descriptions to CWE definitions through
semantic similarity. Third, we construct
smoothed historical CWE exploitation scores
and evaluate their effectiveness in ranking
future vulnerabilities using temporally
separated data.

The main contributions of this paper are:

\begin{itemize}

\item \textbf{Historical CWE exploitability analysis:}
We conduct a large-scale empirical study of
192,788 CVEs, 768 CWE categories, and 3,097
exploited vulnerabilities. We analyze CWE
assignment delays, mapping stability, and
temporal exploitation patterns, and identify
weaknesses significantly associated with
real-world exploitation.

\item \textbf{Hierarchy-aware CWE classification:}
We develop HiCWE, a contrastively trained
bi-encoder that incorporates CWE hierarchical
relationships to improve automated CVE-to-CWE
mapping. On temporally separated test data,
HiCWE achieves 75.3\% top-1 accuracy in flat
CWE classification, outperforming a published
classifier evaluated on the same dataset.

\item \textbf{CWE-informed vulnerability
prioritization:}
We introduce a smoothed historical CWE
exploitability score and evaluate its ability
to prioritize future vulnerabilities.
Using official NVD CWE assignments and rolling
temporal evaluations from 2021 to 2025,
the method achieves ROC-AUC values between
0.692 and 0.727, demonstrating the predictive
value of historical weakness-level information.

\end{itemize}

Together, these contributions establish a
foundation for combining automated root-cause
classification with historical weakness-level
exploitability to support vulnerability
prioritization immediately after disclosure. To support reproducibility and facilitate future research,
we publicly release the experimental code, evaluation
results, and the HiCWE model at
\url{https://anonymous.4open.science/r/cwe-classifier-337D/}.

\section{Overview of the Proposed Framework}
\label{sec:overview}

The objective of this study is to investigate whether
historical exploitation patterns associated with
software weaknesses can support the early prioritization
of newly disclosed vulnerabilities. Unlike conventional
exploitability prediction approaches that primarily
rely on characteristics and threat intelligence
associated with individual CVEs, our approach leverages
the historical exploitation behavior of their underlying
CWE categories. The central hypothesis is that
vulnerabilities sharing a common weakness exhibit
partially consistent exploitation patterns, allowing
historical CWE-level information to provide an early
indication of exploitation risk.

To investigate this hypothesis, we propose a framework
consisting of three interconnected phases, as illustrated
in Figure~\ref{fig:framework_overview}.


\begin{figure*}[!t]
    \centering
    \includegraphics[width=\textwidth]{Figures/overview1.png}
    \caption{Overview of the proposed CWE-informed vulnerability
    prioritization framework.}
    \label{fig:framework_overview}
\end{figure*}


\textbf{Phase 1: Historical CWE Exploitability Analysis.}
We first conduct a large-scale empirical analysis of
CVE--CWE mappings and exploitation evidence collected
from NVD, CISA KEV, VulnCheck, and GreyNoise.
This analysis examines the availability and stability
of CWE assignments, the temporal consistency of
weakness-level exploitation patterns, and the
association between individual CWE categories and
observed exploitation in the wild. Based on these
findings, we construct historical exploitability
profiles that characterize the exploitation behavior
of software weaknesses.

\textbf{Phase 2: Automated CVE to CWE Classification.} Because newly disclosed vulnerabilities may not have an assigned CWE, the second phase infers the underlying weakness directly from the CVE description. The model is based on \texttt{all-MiniLM-L6-v2} and finetuned with Low Rank Adaptation (LoRA). It represents CVE descriptions and CWE definitions in a shared embedding space and uses semantic similarity to identify candidate weaknesses. By leveraging the CWE taxonomy, it first identifies relevant high level categories and then refines the prediction toward more specific CWE identifiers. This enables automated weakness identification without requiring an official CWE assignment.

\textbf{Phase 3: CWE-Informed Exploitability
Prioritization.}
In the final phase, the predicted CWE is linked to
its historical exploitability profile to estimate
the exploitation risk of the newly disclosed CVE.
Historical exploitation rates are estimated from
previously disclosed vulnerabilities using a
smoothed probability calculation to reduce the
influence of sparsely represented weaknesses.
These probabilities are used to rank newly disclosed
CVEs according to their estimated exploitation risk.
The framework is evaluated using temporally separated
data to examine whether historical weakness-level
information can support prioritization of future
vulnerabilities.

Together, these phases connect historical
weakness-level exploitation knowledge with automated
vulnerability classification and early exploitability
assessment. The empirical analysis establishes the
motivation for using CWE history, the bi-encoder
enables the identification of weaknesses for newly
disclosed CVEs, and the prioritization component
translates this information into actionable
vulnerability rankings.

\section{Historical CWE Exploitability Study}
\subsection{Dataset and Methodology}
\label{sec:dataset_methodology}

This empirical study combines vulnerability information from multiple publicly available sources to investigate the evolution of CWE mappings, vulnerability severity, and historical exploitability patterns. The primary source of vulnerability metadata is the National Vulnerability Database (NVD), from which we collected CVE records, publication dates, CWE mappings, CVSS v2/v3/v4 scores, and CVE change-history information. To study the evolution of CWE assignments after disclosure, we additionally used the NVD CVE Change History API, which records modifications made to CVE metadata, including the addition and revision of CWE mappings.

To identify vulnerabilities exploited in the wild, we integrated three complementary exploitation datasets: the CISA Known Exploited Vulnerabilities (KEV) catalog\footnote{https://www.cisa.gov/known-exploited-vulnerabilities-catalog}, VulnCheck's exploitation dataset\footnote{https://research.vulncheck.com/}, and  a proprietary dataset of actively exploited CVEs created by the firm GreyNoise\footnote{ https://www.greynoise.io}. A CVE was considered exploited if it appears in at least one of these sources. After merging the datasets and removing duplicate CVE--CWE pairs, the final dataset contained 192,788 unique CVEs, 768 distinct CWE categories, and 3,097 CVEs known to have been exploited in the wild. These exploited CVEs were used to construct historical exploitability profiles for individual CWE categories.

Different research questions require different temporal subsets of the data. We analyze CVEs published in 2025, in Findings 1 and 2, to minimize the impact of the well-documented NVD processing backlog during 2024, which delayed vulnerability analysis and metadata publication. We examined the complete set of available CVE--CWE mappings in Finding 3, to characterize the distribution of weakness assignments. Finding 4 is based on analysis over temporal exploitability trends using CVEs published between 2020 and 2026, while for Finding 5, we used the complete merged dataset to identify weakness categories that are disproportionately associated with real-world exploitation. 

Prior to analysis, duplicate CVE--CWE pairs were removed, and generic or non-informative CWE labels (\texttt{NVD-CWE-noinfo} and \texttt{NVD-CWE-Other}) were excluded where appropriate. Because a single CVE may be associated with multiple CWE identifiers, analyses involving weakness-level exploitability were performed at the CVE--CWE pair level, whereas analyses involving publication dates, CWE assignment delays, and CVSS scores were performed at the CVE level.


\textbf{CWE Association with Exploitation-} To identify software weakness categories that are disproportionately associated with real-world exploitation, we combined CVE--CWE mappings from the National Vulnerability Database (NVD) with exploitation evidence from CISA Known Exploited Vulnerabilities (KEV), VulnCheck, and GreyNoise. A CVE was considered exploited if it appeared in at least one of these exploitation sources. For each CWE, we computed its exploitation rate, contribution to exploited CVEs, and odds ratio measuring the strength of association with exploitation. Statistical significance was assessed using a one-sided Fisher's exact test, and $p$-values were adjusted using the Benjamini--Hochberg procedure to control the false discovery rate. To ensure robust estimates, only CWEs with at least 100 associated CVEs, at least 10 exploited CVEs, observations spanning at least three years, an odds ratio greater than one, and an adjusted $p$-value below 0.05 were retained. We further evaluated the temporal consistency of exploitation using the proportion of years in which a CWE was exploited and a stability score derived from the coefficient of variation of annual exploitation rates. Eligible CWEs were then ranked using a composite score combining exploitation rate (30\%), odds ratio (25\%), exploited-CVE count (25\%), temporal stability (10\%), and the percentage of years with observed exploitation (10\%). Finally, we compared the resulting ranking with the MITRE 2025 CWE Top 25 using set overlap, Jaccard similarity, and Spearman rank correlation, and performed sensitivity analyses under alternative minimum sample-size thresholds to evaluate the robustness of the findings.

Throughout the empirical study, descriptive statistics are reported using counts, percentages, means, medians, quartiles, and standard deviations where appropriate. Depending on the research question, statistical analyses include Fisher's exact test to evaluate associations between CWE categories and exploitation, Spearman's rank correlation and linear regression to assess temporal stability, Friedman and Wilcoxon signed-rank tests to compare CVSS versions, and bootstrap confidence intervals where applicable. For analyses involving multiple hypothesis tests, $p$-values were adjusted using the Benjamini--Hochberg procedure to control the false discovery rate.

\subsection{Finding 1: Initial CWE Assignments Have Become Faster, but Substantial Delays Persist for Some CVEs}

To motivate the need for automated early CWE assignment, we first examine the delay between the public disclosure of a CVE and the availability of its corresponding CWE mapping. Indeed,  vulnerability prioritization often begins immediately after disclosure. However, CVEs are often published without an assigned CWE, a delay that hinders practitioners from leveraging weakness-level information during the critical early stages of vulnerability assessment. Quantifying these delays provides insight into the potential benefit of automatically predicting the underlying CWE using an LLM at disclosure time.

Figure~\ref{fig:cwe_assignment_delay} Shows the distribution of delays for all CVEs, CVEs published between 2015--2025, and CVEs published in 2025. The distributions are highly right-skewed, with most vulnerabilities receiving a CWE assignment within a few days of publication. Across all CVEs, the median delay was 2.24 days, increasing slightly to 3.17 days for CVEs published between 2015 and 2025. In contrast, the median delay for 2025 CVEs was 0 days, indicating that at least half of newly disclosed vulnerabilities received a CWE assignment on the day of publication. Likewise, the upper quartile decreased from approximately eight days in the historical datasets to essentially zero days in 2025, demonstrating substantial improvements in the timeliness of CWE assignment.

Despite these improvements, a small number of vulnerabilities experience exceptionally long assignment delays, producing mean delays of 65.78, 24.05, and 12.81 days for the three datasets, respectively. These long-tail cases substantially inflate the averages and indicate that weakness-level information may remain unavailable for extended periods for some vulnerabilities. Because CWE is one of the features used by exploitability prediction systems such as EPSS, delays in obtaining an official CWE assignment postpone the availability of weakness-level information. Our LLM-based approach addresses this limitation by inferring the underlying CWE immediately after CVE disclosure, enabling earlier vulnerability prioritization using historical weakness-level exploitability information.
\begin{figure}[!t]
    \centering
    \includegraphics[width=\columnwidth]{Figures/figure 5.jpg}
    \caption{Distribution of delays between CVE publication
    and initial CWE assignment across three time periods.}
    \label{fig:cwe_assignment_delay}
\end{figure}


\subsection{Finding 2: CWE Assignments Are Rarely Revised, and Most Revisions Occur Shortly After Initial Assignment}

Among the 46,600 CVEs published in 2025, only 905 CVEs (around 2\%) experienced at least one change to their CWE mapping after the initial assignment. These updates accounted for 1,006 CWE change events, indicating that a small number of CVEs underwent multiple mapping revisions. Overall, the time lag between the initial CWE mapping and subsequent updates was very short, with a mean of 0.91 days and a median of 0 days, suggesting that most modifications occurred on the same day as the initial mapping. The distribution of update delays was highly concentrated in the first week: 987 (98.1\%) of all CWE change events occurred within 0–7 days, while only 7 (0.7\%) occurred between 8–30 days and 12 (1.2\%) occurred between 31–90 days. No CWE mapping changes were observed after 90 days. These findings indicate that post-publication changes to CWE mappings are relatively uncommon and, when they do occur, they are typically made shortly after the initial CWE assignment.

We limited our analysis to CVEs published in 2025 to avoid bias introduced by the well-documented NVD processing disruption in 2024. The NVD developed a substantial backlog of vulnerability evaluations starting in February 2024, which caused delays in the assignment and updating of metadata, including CWE mappings. If included in our analysis, these operational delays would have artificially inflated the observed time lag between the initial and subsequent CWE mappings. Restricting the analysis to 2025 CVEs circuments the impact of this backlog and provides a more representative view of normal CWE mapping update behavior ~\cite{NVD2024Backlog,NIST2026NVD}.




\subsection{Finding 3: Most CVEs Are Associated with
a Single CWE, but Missing and Multiple Mappings Persist}

Table~\ref{tab:cwe_mapping_distribution_2020_2025} summarizes the number of CWE mappings assigned to each CVE published between 2020 and 2025. Among the 180,599 CVEs analyzed, the majority (138,343; 76.60\%) were associated with only one CWE, indicating that most vulnerabilities are attributed to a single underlying weakness. An additional 18,763 CVEs (10.39\%) were mapped to two CWEs, while only 844 CVEs (0.47\%) were associated with three CWEs.

A notable proportion of vulnerabilities (22,575; 12.50\%) had no specific CWE mapping, reflecting cases where the underlying weakness had not been identified or was recorded using non-informative classifications such as \texttt{NVD-CWE-noinfo} or \texttt{NVD-CWE-Other}. CVEs associated with four or more CWEs were extremely rare, accounting for fewer than 0.05\% of all vulnerabilities. The maximum number of CWE mappings observed for a single CVE was eight, occurring in only one case.

Overall, these results indicate that while most vulnerabilities can be characterized by a single software weakness, a non-negligible fraction involve multiple underlying weaknesses, suggesting that complex vulnerabilities may span multiple stages of a software flaw or affect several related weakness categories.


% \begin{table}[h]
% \centering
% \caption{Distribution of CVEs by Number of CWEs}
% \begin{tabular}{lrrr}
% \toprule
% \textbf{Number of Assigned CWEs} & \textbf{Number of CVEs} \\
% \midrule
% 1   & 282{,}798    \\
% 2   & 36{,}404       \\
% 3   & 1630        \\
% 4   & 134      \\
% 5   & 20        \\
% 6   & 7        \\
% 7   & 6       \\
% 8   & 1        \\
% \bottomrule
% \end{tabular}
% \end{table}


\begin{table}[t]
\centering
\scriptsize
\setlength{\tabcolsep}{3pt}
\caption{Distribution of the Number of CWE Mappings per CVE from 2020--2025}
\label{tab:cwe_mapping_distribution_2020_2025}
\begin{tabular}{lrr}
\toprule
\textbf{Number of Mapped CWEs } & \textbf{  Number of CVEs} & \textbf{  Percentage(\%)} \\
\midrule
0 & 22,575  & 12.50 \\
1 & 138,343 & 76.60 \\
2 & 18,763  & 10.39 \\
3 & 844      & 0.47 \\
4 & 57       & 0.03 \\
5 & 10       & 0.01 \\
6 & 3        & 0.00 \\
7 & 3        & 0.00 \\
8 & 1        & 0.00 \\
\midrule
\textbf{Total} & \textbf{180,599} & \textbf{100.00} \\
\bottomrule
\end{tabular}
\end{table}


% \begin{table}[H]
% \centering
% \caption{Distribution of the Number of CWE Mappings per CVEs in 2025}
% \label{tab:cwe_distribution}
% \begin{tabular}{lrr}
% \toprule
% \textbf{Number of CWE Mappings} & \textbf{Number of CVEs} & \textbf{Percentage (\%)} \\
% \midrule
% 0 & 3,528  & 7.85  \\
% 1 & 34,073 & 75.85 \\
% 2 & 6,904  & 15.37 \\
% 3 & 392    & 0.87  \\
% 4 & 18     & 0.04  \\
% 5 & 2      & 0.00  \\
% 6 & 2      & 0.00  \\
% 7 & 1      & 0.00  \\
% 8 & 1      & 0.00  \\
% \bottomrule
% \end{tabular}
% \end{table}




\subsection{Finding 4: Historical CWE Exploitability Exhibits Moderate Temporal Stability}

To evaluate the temporal stability of CWE exploitability profiles, we analyzed CVEs published between 2020 and 2026 and linked them with known exploitation-in-the-wild records from CISA KEV, VulnCheck, and Greynoise's exploitation dataset. To reduce instability caused by sparse observations, we restricted the analysis to CWE-year combinations containing at least 20 CVEs. This resulted in 908 CWE-year observations spanning 210 unique CWEs.

We used spearman rank correlation to compare CWE exploitability rates between consecutive years. All six year-to-year comparisons exhibited statistically significant positive correlations ($p<0.001$), with correlation coefficients ranging from 0.388 (2022--2023) to 0.568 (2021--2022). Across all consecutive years, the mean Spearman correlation was 0.480, indicating moderate temporal stability in CWE exploitability profiles.

To determine whether historical exploitability predicts future exploitability, we compared exploitability rates in year $t$ with those in year $t+1$. The overall lagged Spearman correlation was 0.468 ($p<0.001$), with a bootstrap 95\% confidence interval of [0.396, 0.538], indicating a consistent positive relationship between historical and future exploitability. Linear regression further confirmed this relationship, producing a statistically significant positive slope of 0.364 ($p<0.001$) and an $R^2$ value of 0.184. Although historical exploitability explains only part of the variation in future exploitability, the results demonstrate that CWEs exhibiting higher exploitation rates in one year tend to remain relatively more exploitable in subsequent years.

Finally, we repeated the analysis using only the 76 CWEs that appeared in every year of the study period. The resulting Spearman correlations ranged from 0.450 to 0.641, consistently exceeding those obtained from the complete dataset. This finding suggests that exploitability profiles are even more stable for well-established CWEs that are consistently represented across years.

Overall, these results indicate that CWE exploitability profiles exhibit moderate but statistically significant temporal stability. Historical exploitation patterns therefore provide useful predictive information for estimating the future exploitability of vulnerability classes, although additional factors beyond CWE also influence whether individual vulnerabilities are exploited in the wild.



\subsection{Finding 5: A Small Set of CWEs Exhibits Persistent Exploitation Associations That Differ from the MITRE Top 25}

\textbf{CWEs Associated with exploitation in the wild-} Among the 192,788 CVEs analyzed, 3,097 were identified as exploited in the wild. To reduce the influence of sparse and potentially unstable estimates, we retained only CWEs associated with at least 100 CVEs and at least 10 exploited CVEs. We further required a positive association with exploitation ($OR>1$) and a Benjamini--Hochberg-adjusted $p$-value below 0.05. Nineteen CWEs satisfied these robustness criteria.

CWE-78 (\emph{Improper Neutralization of Special Elements used in an OS Command}) ranked first according to the composite exploitation score. Of the 4,765 CVEs mapped to CWE-78, 310 were known to be exploited, corresponding to an exploitation rate of 6.51\% and an odds ratio of 4.62. CWE-78 alone accounted for 10.01\% of all exploited CVEs in the dataset and exhibited exploitation in every year of the seven-year study period. CWE-288 (\emph{Authentication Bypass Using an Alternate Path or Channel}) had the highest exploitation rate among the robust categories at 7.87\% and an odds ratio of 5.29, although it accounted for only 1.45\% of all exploited CVEs. Other highly ranked CWEs included CWE-306 (\emph{Missing Authentication for Critical Function}), CWE-434 (\emph{Unrestricted Upload of File with Dangerous Type}), CWE-502 (\emph{Deserialization of Untrusted Data}), CWE-22 (\emph{Path Traversal}), and CWE-287 (\emph{Improper Authentication}).

Several categories combined statistically elevated exploitation likelihood with substantial real-world contribution. CWE-22 (\emph{Path Traversal}) accounted for 6.97\% of exploited CVEs, followed by CWE-862 (\emph{Missing Authorization}) at 5.30\%, CWE-434 (\emph{Unrestricted Upload of File with Dangerous Type}) at 4.84\%, CWE-94 (\emph{Improper Control of Generation of Code ('Code Injection')}) at 4.52\%, CWE-502 (\emph{Deserialization of Untrusted Data}) at 3.78\%, CWE-77 (\emph{Improper Neutralization of Special Elements used in a Command}) at 3.68\%, and CWE-306 (\emph{Missing Authentication for Critical Function}) at 3.52\%. These CWEs were linked to exploitation in each of the seven years of the study, indicating that their observed correlations is not limited to isolated annual events.

Some categories, including CWE-917 (\emph{Improper Neutralization of Special Elements used in an Expression Language Statement ('Expression Language Injection')}), CWE-95 (\emph{Improper Neutralization of Directives in Dynamically Evaluated Code ('Eval Injection')}), and CWE-425 (\emph{Direct Request ('Forced Browsing')}), exhibited relatively high exploitation rates and odds ratios but substantially lower temporal consistency. These CWEs were associated with exploited vulnerabilities in only three or four of the seven study years and had comparatively high coefficients of variation. Therefore, their high aggregate exploitation rates should be interpreted more cautiously than those of CWEs exhibiting persistent exploitation throughout the study period.

\textbf{Comparison with the MITRE CWE Top 25-} Nine of the 19 robust exploitation-associated CWEs also appeared in the MITRE 2025 CWE Top 25: CWE-22 (\emph{Path Traversal}), CWE-77 (\emph{Improper Neutralization of Special Elements used in a Command}), CWE-78 (\emph{Improper Neutralization of Special Elements used in an OS Command}), CWE-94 (\emph{Improper Control of Generation of Code ('Code Injection')}), CWE-306 (\emph{Missing Authentication for Critical Function}), CWE-434 (\emph{Unrestricted Upload of File with Dangerous Type}), CWE-502 (\emph{Deserialization of Untrusted Data}), CWE-862 (\emph{Missing Authorization}), and CWE-918 (\emph{Server-Side Request Forgery (SSRF)}). These shared categories represented 36\% of the MITRE Top 25 and 47.37\% of the robust exploitation-associated set. The Jaccard similarity between the two sets was 0.257. However, the ranks of the nine shared CWEs were not significantly correlated (Spearman's $\rho=-0.067$, $p=0.865$).

The limited ranking agreement reflects the different objectives of the two approaches. The MITRE Top 25 combines weakness prevalence and severity, whereas our ranking focuses specifically on observed exploitation in the wild and incorporates exploitation rate, odds ratio, practical contribution, and temporal consistency. Consequently, some highly prevalent or severe CWEs in the MITRE ranking were not disproportionately associated with observed exploitation, while several CWEs outside the MITRE Top 25 exhibited elevated and persistent exploitation likelihood.

\textbf{Sensitivity Analysis-} We evaluated whether the identified CWE set was sensitive to the minimum sample-size thresholds. Under a relaxed threshold of at least 50 CVEs and five exploited CVEs, all 19 CWEs from the primary analysis remained in the resulting ranking. Increasing the threshold to at least 200 CVEs and 10 exploited CVEs retained 16 of the 19 CWEs (84.21\%), while the strictest threshold of at least 500 CVEs and 20 exploited CVEs retained 15 CWEs (78.95\%). These results indicate that the majority of the identified weakness categories remain associated with exploitation under substantially stricter inclusion requirements, supporting the robustness of the findings.

This finding suggest that vulnerability prioritization should consider historical exploitation characteristics in addition to weakness prevalence and severity. In particular, weaknesses such as CWE-78 (\emph{Improper Neutralization of Special Elements used in an OS Command}), CWE-306 (\emph{Missing Authentication for Critical Function}), CWE-434 (\emph{Unrestricted Upload of File with Dangerous Type}), and CWE-502 (\emph{Deserialization of Untrusted Data}) consistently exhibit elevated exploitation risk across multiple years. Their persistent association with real-world exploitation indicates that these weakness categories should receive high priority during secure software development, vulnerability remediation, and defensive monitoring.




% \begin{table}[H]
% \centering
% \caption{Top 25 CWEs by Mapped Known Exploited Vulnerabilities}
% \begin{tabular}{lrrr}
% \toprule
% \textbf{CWE ID} & \textbf{Mapped KEVs} \\
% \midrule
% 20   & 118    \\
% 78   & 104       \\
% 787   & 100        \\
% 416   & 92      \\
% 119   & 84        \\
% 22   & 75        \\
% 502   & 66       \\
% 94   & 65        \\
% 287   & 39    \\
% 843   & 37       \\
% 306   & 36        \\
% 284   & 33      \\
% 79   & 33        \\
% 264   & 31        \\
% 122   & 30       \\
% 77   & 29        \\
% 89   & 28    \\
% 200   & 24       \\
% 434   & 20        \\
% 918   & 18      \\
% 399   & 18        \\
% 125   & 17        \\
% 288   & 17       \\
% 190   & 16        \\
% 59   & 16        \\
% \bottomrule
% \end{tabular}
% \end{table}

% \begin{table}[H]
% \centering
% \caption{Top 25 CWEs by Mapped Known Exploited Vulnerabilities 2025}
% \begin{tabular}{lrrr}
% \toprule
% \textbf{CWE ID} & \textbf{Mapped KEVs} \\
% \midrule
% 78   & 18    \\
% 502   & 15       \\
% 22   & 13        \\
% 416   & 11      \\
% 787   & 11        \\
% 77   & 7        \\
% 79   & 7       \\
% 287   & 6        \\
% 288   & 6    \\
% 306   & 6       \\
% 94   & 6        \\
% 122   & 5      \\
% 125   & 5        \\
% 121   & 4        \\
% 20   & 4       \\
% 284   & 4        \\
% 36   & 4    \\
% 434   & 4       \\
% 843   & 4        \\
% 862   & 4      \\
% 863   & 4        \\
% 89   & 4        \\
% 119   & 3       \\
% 120   & 3        \\
% 506   & 3        \\
% \bottomrule
% \end{tabular}
% \end{table}



% % \textbf{RQ6}: Can historical CWE exploitability improve the prioritization of newly disclosed CVEs? 
% % This research question evaluates whether incorporating historical exploitability information at the CWE level provides better prioritization of newly disclosed CVEs than using severity-based metrics alone.

% % Methodology: Compute the historical exploitability score for a CWE using only CVEs disclosed before the target CVE. Evaluate CVEs disclosed in 2025 and 2024, its CWE score should be based only on cves before their release date. Extract their CWE exploitability score and check if they really exploitated in the wild after their release. Check if there is any meaningful correlation.

% % if cve-x matched to CWE-y, and for all cves released before the date of cve-x and matched to CWE-y, get their exploitation in the wild. we have exploitation for cve-x in the wild as well. 

% % high correlation or any statistic? 
% % statistical test?

% RQ6: Does incorporating historical CWE exploitability improve the ranking of newly disclosed CVEs for exploitation prioritization compared with CVSS severity alone?








% \subsection{\textbf{RQ6}: How often do vulnerabilities receive lower CVSS scores under newer CVSS versions, and what is the magnitude of these score reductions?}

% As shown in Table~\ref{tab:cvss_score_reduction}, vulnerabilities are substantially more likely to receive lower severity scores under CVSS v4 than under CVSS v3. Among the 5,919 CVEs with both CVSS v2 and v3 scores, 1,111 (18.77\%) received a lower score under CVSS v3. These reductions were generally small, with a mean decrease of 0.39 points, a median of 0.2 points, and a maximum reduction of 2.5 points. In contrast, among the 10,509 CVEs with both CVSS v3 and v4 scores, 7,963 (75.77\%) received a lower score under CVSS v4. The magnitude of these reductions was much larger, with a mean decrease of 2.83 points, a median of 2.9 points, and a maximum reduction of 9.8 points. Overall, the findings indicate that score decreases are relatively uncommon and modest when moving from CVSS v2 to v3, but highly frequent and substantial when moving from CVSS v3 to v4.

% \begin{table}[t]
% \centering
% \caption{Comparison of CVSS Score Reductions Across Versions for CVEs Published in 2025}
% \label{tab:cvss_score_reduction}
% \begin{tabular}{lcc}
% \toprule
% \textbf{Comparison} & \textbf{\% Lower in New Version} & \textbf{Mean Reduction} \\
% \midrule
% CVSS v3 vs. CVSS v2 & 18.77\% & 0.39 \\
% CVSS v4 vs. CVSS v3 & 75.77\% & 2.83 \\
% \bottomrule
% \end{tabular}
% \end{table}

% The substantial reduction in CVSS v4 scores is consistent with the design objectives of the CVSS v4.0 specification. Unlike CVSS v3.1, CVSS v4 introduces the Attack Requirements (AT) metric, refines the calculation of vulnerability impacts, and adopts a redesigned scoring formula that more precisely models exploitability and impact. The specification also separates intrinsic vulnerability severity from contextual factors through enhanced Threat, Environmental, and Supplemental metrics, reducing the reliance on the Base Score alone for prioritization. Collectively, these changes reduce score inflation observed in CVSS v3.1 and produce a broader distribution of Base Scores, explaining why a large proportion of vulnerabilities receive lower severity scores under CVSS v4.0.



% % \subsection{\textbf{RQ7}: How significant are the differences in vulnerability severity scores across CVSS versions?}

% % A Friedman test was conducted to evaluate whether vulnerability severity scores differ across CVSS versions. The results revealed a statistically significant difference among CVSS v2, v3, and v4 scores ($\chi^2 = 6379.49$, $p < 0.001$), indicating that the three versions do not assign equivalent severity scores. Pairwise Wilcoxon signed-rank tests with Bonferroni correction further confirmed statistically significant differences between both version pairs ($p < 0.001$), as summarized in Table~\ref{tab:cvss_version_comparison}. Compared with CVSS v2, CVSS v3 assigned higher severity scores to 4,773 (80.64\%) of the 5,919 comparable vulnerabilities, lower scores to 1,111 (18.77\%), and identical scores to 35 (0.59\%). The mean severity score increased from 6.19 to 7.83, corresponding to a large positive effect size (rank-biserial correlation $r = 0.903$; Cohen's $d = 1.246$). In contrast, compared with CVSS v3, CVSS v4 assigned lower severity scores to 7,963 (75.77\%) of the 10,509 comparable vulnerabilities, higher scores to 2,431 (23.13\%), and identical scores to 115 (1.09\%). The mean severity score decreased from 7.51 to 5.61, with a large negative effect size (rank-biserial correlation $r = -0.714$; Cohen's $d = -0.773$). These findings indicate that CVSS v3 systematically increased vulnerability severity relative to CVSS v2, whereas CVSS v4 systematically reduced severity scores relative to CVSS v3.

% % \begin{table}[t]
% % \centering
% % \caption{Comparison of Vulnerability Severity Scores Across CVSS Versions}
% % \label{tab:cvss_version_comparison}
% % \begin{tabular}{lrrrrrr}
% % \toprule
% % \textbf{Comparison} & \textbf{N} & \textbf{\% Higher} & \textbf{\% Lower} & \textbf{$r$} & \textbf{$d$} & \textbf{$p$-value} \\
% % \midrule
% % CVSS v2 vs. CVSS v3 &
% % 5,919 &
% % 80.64 &
% % 18.77 &
% % 0.903 &
% % 1.246 &
% % $<0.001$ \\

% % CVSS v3 vs. CVSS v4 &
% % 10,509 &
% % 23.13 &
% % 75.77 &
% % $-0.714$ &
% % $-0.773$ &
% % $<0.001$ \\
% % \bottomrule
% % \end{tabular}
% % \end{table}





% \subsection{Discussion}


% Our findings in RQ4 provide empirical justification for incorporating historical CWE exploitability into CVE risk assessment. Although exploitability is not completely determined by the weakness type, CWE exploitability profiles exhibit moderate and statistically significant temporal stability across multiple years. Furthermore, historical exploitability rates are positively correlated with future exploitability ($\rho=0.468$, $p<0.001$), indicating that CWEs with consistently higher exploitation rates tend to remain relatively more attractive to attackers over time. These results suggest that the historical exploitability of a CWE constitutes a meaningful predictive feature for estimating the likelihood that newly disclosed CVEs associated with that CWE will be exploited in the wild. Nevertheless, the observed $R^2$ value of 0.184 indicates that CWE history alone is insufficient for accurate prediction and should be combined with additional vulnerability characteristics, such as CVSS metrics, affected products, and exploit availability, to improve predictive performance.



\subsection{Discussion}
\label{sec:empirical_discussion}

The empirical findings provide the motivation and
supporting evidence for incorporating historical
CWE exploitability into early vulnerability
prioritization. Our analysis highlights two key
aspects of this approach: the availability and
reliability of CWE assignments, and the extent
to which historical exploitation patterns can
inform the assessment of newly disclosed
vulnerabilities.

First, although the timeliness of initial CWE
assignments has improved, some vulnerabilities
continue to experience substantial delays.
While the median assignment delay for CVEs
published in 2025 was zero days, the mean
remained 12.81 days, reflecting a long-tailed
distribution. Consequently, weakness-level
information may not be available when a
vulnerability is first disclosed. This finding
motivates automated CWE prediction to enable
the use of weakness-level information before
an official assignment becomes available.

Second, CWE mappings are generally stable after their initial assignment. Only 2.01\% of CVEs published in 2025 experienced subsequent CWE mapping changes, with most revisions occurring within the first week. Furthermore, 76.60\% of CVEs published between 2020 and 2025 were associated with exactly one CWE. These findings suggest that single-label CWE prediction is a reasonable starting point for automated mapping, although missing and multiple CWE assignments must also be considered. Importantly, the stability of official assignments does not guarantee the accuracy of automated predictions, which must be evaluated independently.

More importantly, our analysis demonstrates
that historical CWE exploitability exhibits
moderate temporal stability. Exploitation rates
were positively correlated across consecutive
years, with an overall lagged Spearman correlation
of $\rho=0.468$ ($p<0.001$). This indicates
that CWE categories with relatively high
exploitation rates tend to maintain higher
rates in subsequent years. In addition, we
identified 19 CWE categories with statistically
significant and robust associations with
exploitation in the wild. Several of these
weaknesses, including command injection,
authentication-related weaknesses, unrestricted
file upload, deserialization, and path traversal,
exhibited persistent exploitation across
multiple years.

The limited agreement between our exploitation-
based CWE ranking and the MITRE 2025 CWE Top 25
further highlights the importance of considering
observed exploitation patterns. While the MITRE
ranking emphasizes weakness prevalence and
severity, our analysis identifies categories
associated with elevated historical exploitation
likelihood. These perspectives are complementary,
but they serve different vulnerability management
objectives.

Taken together, these findings support the use
of historical CWE exploitation rates as an
informative signal for early vulnerability
prioritization. Although historical CWE rates
explain only part of the variation in subsequent
annual exploitation rates ($R^2=0.184$),
their statistically significant temporal
association suggests that they can provide
useful information when vulnerability-specific
exploitation evidence is limited or unavailable.
The objective is therefore not to determine
exploitation solely from the CWE identifier,
but to use historical weakness-level knowledge
to support the initial ranking of newly
disclosed vulnerabilities.

These observations establish the motivation
for the next component of our framework:
automatically identifying the underlying CWE
of a newly disclosed CVE from its description.
By linking the predicted CWE to its historical
exploitability profile, the proposed approach
can assign an initial exploitation-based
priority without waiting for an official CWE
mapping. The subsequent sections investigate
the accuracy of automated CWE identification
and the effectiveness of historical CWE
information for vulnerability prioritization.



\section{Automating CVE to CWE mapping}

The objective of our framework is to automatically map a newly disclosed CVE to its underlying Common Weakness Enumeration (CWE) using a Large Language Model (LLM), thereby enabling exploitability estimation before real-world exploitation evidence becomes available. As demonstrated in the previous section, different CWE categories exhibit substantially different historical exploitation characteristics. Therefore, once a CVE is accurately mapped to its corresponding CWE, the historical exploitability profile associated with that weakness can be used to estimate the likelihood that the new vulnerability will be exploited in the wild. Consequently, the accuracy of CWE classification directly determines the reliability of the proposed exploitability estimation framework.

In our proposed automated CWE mapping approach, we leverage the hierarchical organization of the CWE taxonomy rather than requiring the model to select directly from hundreds of weakness entries. According to the MITRE CWE Usage Guidance~\cite{mitre_cwe_guidance}, weaknesses are organized into multiple levels of abstraction, including \emph{Pillar}, \emph{Class}, \emph{Base}, and \emph{Variant}. Pillar entries represent broad conceptual categories, Class entries group related weaknesses, Base entries describe specific weakness types suitable for vulnerability classification, and Variant entries provide technology- or implementation-specific refinements. This hierarchical structure enables a coarse-to-fine classification strategy that progressively narrows the candidate search space while preserving semantic relationships among weaknesses. MITRE provides several CWE views, each organizing weaknesses according to a particular perspective or application domain. In this study, we adopt \emph{CWE-699: Software Development}, which contains 40 weakness categories and 399 weaknesses. By restricting our analysis to this software-oriented view, we establish a well-defined classification scope, reduce the candidate search space, and facilitate a more controlled and consistent evaluation of CWE mapping performance.





\subsection{HiCWE: Hierarchy-Aware Contrastive CWE Encoder}

In designing our automated CWE mapping framework for newly disclosed CVEs, we considered the following key requirements:

\begin{itemize}
    \item \textbf{Adaptable to new CWEs:} New CWE categories are rare, but they do occur and having a system that is capable of distinguishing new categories without retraining is crucial.
    \item \textbf{Independent of syntax and language:} CVEs often include product names, versions, and irrelevant information that could throw off the system. Therefore, the system should not rely on syntax, hard-coded rules, and keywords during inference.
    
\item \textbf{Generalization and no memorization:}
The system must prevent overlap between training and
test datasets and learn generalizable semantic
relationships between CVE descriptions and CWE
definitions, rather than relying on memorized
CVE--CWE associations.

\item \textbf{Efficient semantic matching:}
    The framework should efficiently match CVE
    descriptions with CWE definitions using semantic
    similarity rather than fixed-label classification.

\end{itemize}


We adopt a bi-encoder approach to efficiently match
CVE descriptions with CWE definitions through semantic
similarity rather than fixed-label classification.
Sentence-BERT~\cite{reimers2019sentence} demonstrated
the effectiveness of bi-encoders for semantic retrieval,
while SimCSE~\cite{gao2021simcse} showed that contrastive
learning improves sentence representations. Building
on these findings, we use a contrastively trained
encoder to distinguish semantically similar CWEs
and accommodate new CWE definitions without retraining
a fixed classification layer.



\subsubsection{Methodology}
We used an embedding based approach to map CVE to CWE descriptions shown in Figure \ref{fig:methodology}. This approach allows us to employ a pipeline similar to RAG, where we use cosine similarity to match descriptions together; however, \textit{we do not use} LLMs to generate the final match. I
In this paper, we fine-tune the \texttt{all-MiniLM-L6-v2}
model using CVEs disclosed between 1999 and 2021
and achieve superior CWE classification performance
compared with the evaluated baseline models.

% Methodology figure
\begin{figure}[H]
\centering
\includegraphics[width=1\columnwidth]{Figures/methodology/methodology.png}
\caption{CVE $\rightarrow$ CWE Mapping Methodology}
\label{fig:methodology}
\end{figure}

\subsubsection{CWE Dataset: }

We gather all of our CWE data from \texttt{cwe.mitre.org} under the \texttt{CWE-699} view, \textit{Software Development}. 
We scrape everything including categories and children to visualize their relationship in Figure \ref{fig:cwe_dependency_graph}, which, most importantly, shows that some weaknesses are under multiple categories. Children with multiple categories pose a problem because we cannot assume that each child has a single parent in our multistage search when finding the right CWE for a given CWE description. In order to solve this issue, we use BFS to first choose the parent that the model perceives as most appropriate, then fetch the children associated with that parent. Additionally we filter our train and test dataset to only include examples that have children in CWE-699. % in order to isolate for noise and reduce candidates.

\begin{figure}[H]
\centering
\includegraphics[width=1\columnwidth]{Figures/methodology/cwe_dependency_graph.png}
\caption{CWE-699 Dependency Graph}
\label{fig:cwe_dependency_graph}
\end{figure}


\subsubsection{CVE Dataset}

We use a training set of CVEs disclosed from 1999 to 2021 and a validation set of CVEs disclosed from 2023 to 2025. The wide time span of the validation set lets us evaluate the finetuned model on a large and varied sample. Our training data is based on the dataset of \texttt{zefang-liu}~\cite{liu2024cvecwedataset}, to which we add CWE descriptions and apply further postprocessing to remove as much noise as possible. We use this dataset in part so that our state of the art baseline, described in the model section, is trained on the same data, which keeps the comparison consistent.

\textbf{Post processing:} Vulnerability descriptions often contain product names, vendor names, and version numbers that are strongly correlated with specific CVEs, so a fine-tuned model can learn to "recognize" a vulnerability by memorizing these surface cues rather than by understanding the flaw itself. Such shortcuts would not generalize to new products or unseen vulnerabilities, and they would inflate our evaluation results. To force the model to learn from the technical mechanism instead, we apply a rule-based filter to each description. First, we delete all parenthetical text, which typically carries auxiliary notes such as patch advice. Next, we remove every capitalized word that contains a lowercase letter (e.g., \texttt{Apache}, \texttt{Log4j2}), except for the first word of each sentence and words on a whitelist of vulnerability-class terms (e.g., \emph{Injection}, \emph{Overflow}), common CVE boilerplate, and standard technology acronyms (e.g., SQL, HTTP, JSON, LDAP), which describe the weakness rather than identify the product. We then remove any alphanumeric token containing a digit, which strips version numbers and names such as \texttt{log4j-core}, and finally collapse leftover whitespace. In the Log4Shell example, this removes the product name, version ranges, and parenthetical, while preserving the description of the flaw itself.

\subsubsection{Training Method}

We trained on 60,092 CVEs disclosed between 1999 and 2021. The corresponding taxonomy contains 146 CWEs grouped under 35 categories, for a total of 181 nodes, and 25 of those nodes belong to more than one category. We split the CVEs by identifier so that a CVE with several labels never appears in two splits, and we stratified the split on the primary CWE so that rare weaknesses are still present in validation and test. This leaves 48,066 CVEs for training, 6,013 for validation, and 6,013 for testing. Every node is rendered as a short text that joins its title, its summary, and the titles of its immediate children, and all 181 node texts are distinct. We train with a contrastive loss over each batch: each CVE is pulled toward its own node and pushed away from the other nodes in the batch. Because those other nodes act as negatives, we keep the node texts distinct so that no negative accidentally equals the positive. We also mined harder negatives with the untuned encoder, taking for each CWE the closest nodes that are neither related nor siblings, which look similar but belong to another category, along with the closest unrelated categories.

From these pairs we built two training recipes. Children only uses plain CVE to own CWE pairs and relies on the batch for negatives, giving 48,686 pairs. Hierarchy full keeps those pairs and adds category level data: 48,686 CVE to child pairs, each paired with a sibling and a look alike from another category as negatives; 50,497 CVE to category pairs, each contrasted against the nearest unrelated category; and 404 CWE to sibling pairs that pull siblings together and push look alikes apart. Both recipes share the split, the base model, the optimizer, and the seed, so the comparison isolates the effect of the training pairs.

\subsubsection{Model}

\textbf{Baseline:} We evaluate against two baselines. The first is the untuned \texttt{all-MiniLM-L6-v2} and \texttt{all-MiniLM-L12-v2} encoders under the same retrieval procedure. We chose MiniLM because it is a small distilled bi-encoder that is strong at semantic similarity for its size, cheap to run, and offered in two sizes in one family, which lets us vary capacity without changing anything else; its untuned version therefore isolates the contribution of finetuning. The second baseline is Mulliken's CWE Predictor~\cite{mullikencwepredictor}, a finetuned DistilBERT classifier trained over 232 CWE categories. Since it returns only a flat ranking of children, we map each predicted child to its category to obtain a category ranking.

\textbf{Model finetuning:} We finetune \texttt{all-MiniLM-L6-v2} and its larger sibling \texttt{all-MiniLM-L12-v2} with LoRA to show how the recipe scales with model size. We train on the 1999 to 2021 CVEs set aside for this purpose and merge the adapters into the base weights before evaluation, so inference is a single forward pass through the encoder. The finetuned models are released on Hugging Face~\cite{hicwe2026models}.

\textbf{Results-} Table~\ref{tab:accuracy} \textbf{reports Accuracy@$k$ on CVEs disclosed from 2023 to 2025}, which the models did not see during training. 

% Figure~\ref{fig:recall} shows the corresponding curves.

Finetuning improves every target over the untuned encoder. At $k=1$, the untuned baseline finds few correct categories and children, while both recipes improve on it by a wide margin. This shows that the retrieval signal can be learned from our data.

The flat child target gives the highest accuracy of the child targets. Here the model searches the whole taxonomy directly rather than only the children of the top categories, so it is the most direct test of child prediction. Even on this target, recipe A (hier full), which is trained with the taxonomy structure, finds more correct children than recipe B (children only), which is trained on flat pairs. Training with the hierarchy therefore improves child prediction in general, not only when the search is restricted to a few categories. Recipe A also finds more correct children than the SOTA classifier.

The difference comes from category placement. Recipe A learns to place a CVE near its parent category, while recipe B does not. Recipe B therefore loses much accuracy when the search is restricted to the top categories, and a smaller amount in the flat setting.

On the category target, the SOTA classifier is slightly better at $k=1$ and $k=2$. Recipe A matches it at $k=3$ and surpasses it from $k=5$, where SOTA stops improving.

\textbf{Effect of encoder size: }
Since recipe A gave the largest gains on MiniLM L6, we trained the same recipe on MiniLM L12, which has twice as many layers and a larger capacity, to see how the method scales. The untuned L12 is clearly stronger than the untuned L6 (for example, category Accuracy@1 of 0.39 vs 0.24), but after finetuning the gap almost disappears. For recipe A, the gains from L6 to L12 are at most about 4 points on the category and hierarchical child targets, and the flat child target is essentially unchanged (0.75 at $k=1$ and 0.85 at $k=5$ for both sizes). Doubling the depth therefore adds little once the model is trained with the hierarchy. We hypothesise that the limit is the training data rather than the model capacity.

Overall, using the hierarchy during training gives the strongest child predictions, which is the quantity that matters for the exploitability step.

\begin{table}[H]
\centering

\caption{CVE$\rightarrow$CWE Accuracy@$k$ for CVEs disclosed from 2023 to 2025, which the models did not see during training. Best value in each column in bold. The published classifier has no category stage, so it cannot produce hierarchical child accuracy.}
\label{tab:accuracy}
\begin{tabular}{llccc}
\toprule
 &  & @1 & @3 & @5 \\
\midrule
\multirow[c]{6}{*}{Category} & Baseline (MiniLM L6) & 0.236 & 0.409 & 0.510 \\
 & Baseline (MiniLM L12) & 0.386 & 0.632 & 0.744 \\
 & A: heir full (MiniLM L12) & 0.789 & \textbf{0.867} & \textbf{0.897} \\
 & A: heir full (MiniLM L6) & 0.750 & 0.856 & 0.889 \\
 & B: children only (MiniLM L6) & 0.584 & 0.776 & 0.841 \\
 & SOTA (cwe-predictor) & \textbf{0.802} & 0.862 & 0.869 \\
\cline{1-5}
\multirow[c]{5}{*}{Child (hier.)} & Baseline (MiniLM L6) & 0.124 & 0.175 & 0.216 \\
 & Baseline (MiniLM L12) & 0.236 & 0.306 & 0.344 \\
 & A: heir full (MiniLM L12) & \textbf{0.726} & \textbf{0.770} & \textbf{0.787} \\
 & A: heir full (MiniLM L6) & 0.695 & 0.735 & 0.756 \\
 & B: children only (MiniLM L6) & 0.506 & 0.557 & 0.584 \\
\cline{1-5}
\multirow[c]{6}{*}{Child (flat)} & Baseline (MiniLM L6) & 0.337 & 0.432 & 0.479 \\
 & Baseline (MiniLM L12) & 0.346 & 0.446 & 0.496 \\
 & A: heir full (MiniLM L12) & \textbf{0.754} & 0.824 & 0.850 \\
 & A: heir full (MiniLM L6) & 0.753 & \textbf{0.828} & \textbf{0.853} \\
 & B: children only (MiniLM L6) & 0.715 & 0.795 & 0.823 \\
 & SOTA (cwe-predictor) & 0.655 & 0.753 & 0.779 \\
\cline{1-5}
\bottomrule
\end{tabular}
\end{table}








\subsection{Why Not Use Generative LLMs for CWE Mapping?}
\label{subsec:generative_llm_limitations}


Generative large language models (LLMs), such as ChatGPT,
offer a straightforward approach to automated CVE-to-CWE
mapping: given a vulnerability description, the model
can be prompted to identify the corresponding software
weakness. Previous research has investigated the use of
prompt-based LLMs for vulnerability description mapping,
including CWE identification~\cite{liu2023not}.
However, CWE classification requires selecting a precise
identifier from a large, hierarchical taxonomy containing
semantically overlapping weaknesses. Generative models
may confuse CWE identifiers and definitions, select
overly general weakness categories, or produce
inconsistent predictions. To further investigate whether
prompt-based generative LLMs provide sufficiently reliable
CWE assignments, we conducted a preliminary experimental
evaluation of several locally deployed models.


We evaluated Llama 3.2 (3B), Qwen 2.5 (14B),
Qwen 2.5-Coder (14B), and Phi-4 (14B) using randomly selected CVEs with their NVD CWE assignments
as reference labels. Preliminary experiments showed
that allowing LLMs to generate CWE identifiers
without an explicit candidate list frequently
resulted in confusion between CWE identifiers,
names, and descriptions. We therefore developed
a constrained, hierarchy-aware prompting pipeline
that supplied official CWE identifiers, titles,
and descriptions. 

Configuration A uses a two-stage approach, selecting one
category from 40 CWE-699 categories as parent CWE and then identifying
the most relevant child CWE. Configuration B expands
the parent candidate set to 157 software-related CWEs
by incorporating 116 additional Pillar- and Class-level
weaknesses from CWE-699. Configuration C reduces error
propagation by selecting the top three parent candidates,
combining their children, and performing a final selection
over the resulting candidate set.



The initial evaluation on 25 CVEs revealed limited
classification accuracy for both models. Under
Configuration A, Llama 3.2 (3B) achieved 0\%
exact-match accuracy, while Qwen 2.5 (14B) achieved
4\%. Expanding the taxonomy in Configuration B
improved exact-match accuracy to 12\% and 16\%,
respectively. However, Llama 3.2 repeatedly predicted
the same CWE for 24 of the 25 CVEs, indicating that
its apparent improvement did not reflect reliable
weakness discrimination. Moreover, both configurations
restricted classification to the children of a single
predicted parent, making an incorrect initial selection
difficult to recover from.

To address this limitation, Configuration C retains
the three most relevant parent candidates and performs
a final selection among their combined children.
In the initial evaluation, Qwen 2.5 achieved 20\%
exact-match accuracy under Configuration C, compared
with 16\% under Configuration B. Although Llama 3.2's
accuracy decreased to 8\%, its prediction diversity
improved substantially. Based on these findings,
we selected Configuration C for further evaluation
on a larger dataset of 1000 CVEs to assess its
performance across additional LLM architectures.


For configuration C over 1000 CVEs, at temperature $T=0$, Qwen 2.5-Coder
achieved the highest exact-match accuracy (23.3\%)
and parent-level accuracy (28.0\%), outperforming
Qwen 2.5 (17.3\% and 20.7\%), Phi-4 (16.7\%
and 16.0\%), and Llama 3.2 (7.3\% and 9.3\%).
Adjusting the decoding temperature improved
Qwen 2.5-Coder's exact-match accuracy to 26.0\%
at $T=0.3$. Furthermore, enriching the CVE
description with information retrieved from
external reference links increased exact-match
accuracy to 32.0\% and parent-level accuracy
to 32.7\%, although average processing time
increased from 3.9 to 6.4 seconds per CVE.

Despite these refinements, the best-performing
generative configuration correctly identified
the exact CWE for only 32.0\% of the evaluated
vulnerabilities. Moreover, 10.7\% of the
reference CWE assignments were outside the
evaluated taxonomy, illustrating the importance
of candidate coverage. These findings suggest
that prompt engineering, hierarchical reasoning,
and contextual enrichment alone do not guarantee
reliable fine-grained CWE identification.
The observed sensitivity to model selection,
taxonomy construction, and decoding parameters
also complicates the use of generative models
as a consistent automated classification mechanism.

These limitations motivate our investigation of
a lightweight, fine-tuned bi-encoder as an
alternative to generative CWE prediction.
Rather than generating CWE identifiers through
autoregressive decoding, the proposed approach
learns semantic representations of CVE descriptions
and CWE definitions and ranks predefined candidate
weaknesses according to their similarity.
This design avoids unconstrained identifier
generation and provides a deterministic
candidate-ranking mechanism. The following
section introduces the proposed bi-encoder
architecture and evaluates whether task-specific,
parameter-efficient fine-tuning can improve
CWE identification.


\begin{table}[t]
\centering
\caption{Generative LLM performance for CVE-to-CWE
mapping ($n=1000$). Results use Configuration C
and $T=0$, unless otherwise specified.}
\label{tab:generative_llm_summary}
\scriptsize
\setlength{\tabcolsep}{3pt}
\begin{tabular}{lrr}
\toprule
\textbf{Model / Configuration}
& \textbf{Exact (\%)}
& \textbf{Parent (\%)} \\
\midrule
Llama 3.2 (3B)        & 7.3  & 9.3 \\
Qwen 2.5 (14B)        & 17.3 & 20.7 \\
Phi-4 (14B)          & 16.7 & 16.0 \\
Qwen 2.5-Coder (14B)  & 23.3 & 28.0 \\
\midrule
Qwen 2.5-Coder ($T=0.3$)
                       & 26.0 & 26.7 \\
+ Reference enrichment & 32.0 & 32.7 \\
\bottomrule
\end{tabular}
\end{table}


It is important to note that the evaluated generative
LLMs may have encountered some of the CVE descriptions
and their corresponding CWE assignments during
pretraining, as this information is publicly available
through vulnerability databases. Although the extent
of such exposure cannot be verified, the possibility
of training-data overlap makes the relatively low
classification accuracy particularly noteworthy.
Even with potential prior exposure to some CVE--CWE
mappings, the best-performing configuration achieved
only 32.0\% exact-match accuracy. This observation
further motivates the development of a task-specific
CWE classification approach.

To address this concern in our proposed bi-encoder
evaluation, we adopt a temporal data-splitting
strategy, using CVEs published between 1999 and
2021 for training and CVEs published between 2023
and 2025 for testing. This separation prevents
direct overlap between the CVE records used for
fine-tuning and those used for evaluation, providing
a more realistic assessment of the model's ability
to generalize to newly disclosed vulnerabilities.


The results of the temperature sensitivity experiments
under Configuration C for the four evaluated LLMs are
reported in Tables~\ref{tab:temperature_all}, in the Appendix.
To facilitate reproducibility, all prompts used in
the generative LLM experiments, together with the
implementation and evaluation scripts, are publicly
available in our repository at
\url{https://anonymous.4open.science/r/cwe-classifier-337D/}.




\section{Exploitability Ranking Using Historical CWE Exploitability Score}
\label{sec:exploitability-ranking}

To evaluate the practical usefulness of historical weakness-level exploitability, we conducted a temporal exploitability ranking experiment. Unlike supervised exploitability prediction models, the objective of this experiment is to determine whether the historical exploitability associated with a software weakness (CWE) can directly prioritize newly disclosed vulnerabilities. This experiment also serves as the foundation for evaluating our LLM-based framework, where official CWE assignments can later be replaced with automatically predicted CWE labels.

% For each newly disclosed CVE, the associated CWE is mapped to its historical exploitability probability computed from vulnerabilities published in previous years. CVEs are then ranked according to this historical exploitability score, and the resulting ranking is evaluated against vulnerabilities that were subsequently observed to be exploited in the wild.

In this experiment, we used the official CVE--CWE mappings provided by the National Vulnerability Database (NVD). We considered a vulnerability is exploited if it appears in at least one of the publicly available exploitation datasets used throughout this paper, namely CISA Known Exploited Vulnerabilities (KEV), VulnCheck, or GreyNoise.

To prevent temporal information leakage, historical exploitability is computed using only vulnerabilities published before the evaluation year. Specifically, for each CWE $c$, the historical exploitability probability is estimated as

\begin{equation}
\small
P(\mathrm{Exploited}\mid c)=
\frac{E_c+\alpha P_0}
{N_c+\alpha},
\label{eq:historical_exploitability}
\end{equation}

where $E_c$ denotes the number of historically exploited CVEs mapped to CWE $c$, $N_c$ is the total number of historical CVEs associated with that CWE, $P_0$ is the overall historical exploitation rate across all historical CVEs, and $\alpha$ is a smoothing parameter.


The smoothing parameter prevents unreliable exploitability estimates for weaknesses with only a few historical observations. Without smoothing, a rarely observed CWE could receive an unrealistically high exploitability probability based on only one or two exploited vulnerabilities. The parameter $\alpha$ incorporates prior knowledge by adding pseudo-observations equivalent to the global historical exploitation rate. As the number of historical CVEs associated with a CWE increases, the influence of the smoothing term becomes negligible, and the estimated exploitability converges to the observed historical exploitability of that weakness. In this study, we performed a sensitivity analysis using $\alpha \in \{20,40,60\}$. The results showed that increasing the smoothing parameter from 20 to 40 improved the exploitability ranking performance, while further increasing it to 60 produced negligible changes. Therefore, the remaining experiments report the results obtained using $\alpha=40$.

Each CVE $v$ in the evaluation year is assigned the historical exploitability score of its mapped CWE. When multiple CWE mappings exist for a CVE, the maximum historical exploitability score is used. The ranking score assigned to each vulnerability is therefore defined as

\begin{equation}
\small
\mathrm{RankScore}(v)=
\max_{c \in \mathrm{CWE}(v)}
P(\mathrm{Exploited}\mid c),
\label{eq:ranking_score}
\end{equation}

where $\mathrm{CWE}(v)$ denotes the set of CWE mappings associated with vulnerability $v$. CVEs are then ranked in descending order according to $\mathrm{RankScore}(v)$.


The experiment follows a rolling temporal evaluation strategy. Historical exploitability probabilities are computed using all CVEs published before the evaluation year, while testing is performed exclusively on CVEs published during that year. Specifically, five independent evaluations are performed for the period 2021--2025. For example, when evaluating CVEs published in 2023, historical exploitability probabilities are estimated using only vulnerabilities disclosed before 2023. This temporal evaluation prevents information leakage from future vulnerabilities while assessing the robustness of historical weakness-level exploitability over multiple years.


The ranked vulnerabilities are evaluated using the Area Under the Precision--Recall Curve (PR-AUC), Area Under the ROC Curve (ROC-AUC), Precision@$k$, Recall@$k$, and Normalized Discounted Cumulative Gain (NDCG@$k$), where $k \in \{50,100,250,500\}$. Because exploited vulnerabilities constitute only a small fraction of all disclosed CVEs, PR-AUC is considered the primary evaluation metric. The evaluation cutoffs assess the quality of the ranking produced by Equation~(\ref{eq:ranking_score}) under different remediation capacities. Specifically, Precision@$k$, Recall@$k$, and NDCG@$k$ evaluate the top $k$ highest-ranked vulnerabilities according to their historical exploitability scores. Smaller values of $k$ (50 and 100) measure the effectiveness of highly selective prioritization, whereas larger values (250 and 500) evaluate the ranking quality when organizations can investigate or remediate a broader set of vulnerabilities.


\textbf{Results-} Table~\ref{tab:cwe_ranking} summarizes the exploitability ranking performance obtained using official NVD CWE assignments across five temporal evaluation years (2021--2025). For each evaluation year, historical exploitability probabilities were computed exclusively from vulnerabilities published prior to the corresponding test year, ensuring a leakage-free assessment.

\begin{table*}[h]
\centering
\caption{Temporal exploitability ranking using official NVD CWE assignments.}
\label{tab:cwe_ranking}
\scriptsize
\begin{tabular}{lrrrrrrrr}
\toprule
\textbf{Year} &
\textbf{CVEs} &
\textbf{  Exploited-CVEs} &
\textbf{ PR-AUC} &
\textbf{ ROC-AUC} &
\textbf{ P@100} &
\textbf{ R@100} &
\textbf{ NDCG@100} &
\textbf{ Coverage} \\
\midrule
2021 & 21,019 & 476 & 0.0520 & 0.6919 & 0.050 & 0.0105 & 0.0857 & 79.9\% \\
2022 & 25,793 & 448 & 0.0466 & 0.7016 & 0.050 & 0.0112 & 0.0964 & 84.5\% \\
2023 & 29,377 & 534 & 0.0401 & 0.7036 & 0.030 & 0.0056 & 0.0573 & 88.4\% \\
2024 & 40,484 & 677 & 0.0408 & 0.7161 & 0.000 & 0.0000 & 0.0688 & 93.2\% \\
2025 & 46,600 & 631 & 0.0352 & 0.7268 & 0.110 & 0.0174 & 0.0994 & 90.1\% \\
\bottomrule
\end{tabular}
\end{table*}

Overall, the results demonstrate that historical weakness-level exploitability provides a useful signal for prioritizing newly disclosed vulnerabilities. Across all five evaluation years, the proposed ranking consistently achieved ROC-AUC values between 0.692 and 0.727, indicating that historical exploitability effectively distinguishes vulnerabilities that are later exploited from those that are not.

The Precision--Recall performance ranged from 0.035 to 0.052 across the evaluation period. Although PR-AUC gradually decreased over time, this trend coincides with a substantial reduction in the prevalence of exploited vulnerabilities, from 2.26\% in 2021 to 1.27\% in 2025, making the ranking task increasingly challenging. Relative to the expected PR-AUC of a random ranking, which equals the prevalence of exploited vulnerabilities, the proposed approach consistently achieved approximately 2.2--2.8 times better performance, demonstrating that historical CWE exploitability contains meaningful information for identifying vulnerabilities that are eventually exploited in the wild.

The ranking quality also remained stable across different evaluation cutoffs. For example, in 2025, 11 of the top 100 ranked vulnerabilities (Precision@100 = 0.11) were subsequently observed to be exploited, despite an overall exploitation prevalence of only 1.27\%. Similarly, the NDCG@100 values ranged from 0.057 to 0.099 across the five evaluation years, indicating that exploited vulnerabilities were consistently concentrated near the top of the ranked lists.

Finally, the applicability of the approach improved over time as the availability of CWE mappings increased. The percentage of CVEs with usable CWE assignments grew from approximately 80\% in 2021 to over 90\% in 2025, demonstrating that historical weakness-level exploitability can be applied to the majority of newly disclosed vulnerabilities.

These results establish an upper bound on exploitability ranking performance when official NVD CWE assignments are available. 
In the next section, we replace the official CWE mappings with the CWE labels predicted by our hierarchy-aware LLM framework to evaluate whether automated CWE classification can provide comparable exploitability ranking performance immediately after CVE disclosure.


%%%%%%%%%%%%%%%%%%%%%%%%%%%%%%%%%%
\textbf{Comparison with EPSS-} We further investigated the numerical agreement between historical CWE-based exploitability scores and EPSS scores. For each CVE,
we calculated the absolute difference between the two scores,
expressed in percentage points. CVEs with zero historical
CWE-based scores were excluded from the comparison. The absolute difference between the historical CWE-based exploitability score ($\mathrm{RankScore}(v)$) and the EPSS score ($\mathrm{EPSS}(v)$) is measured in percentage points (pp), calculated as $D(v)=100\times|\mathrm{RankScore}(v)-\mathrm{EPSS}(v)|$. For example, scores of 0.08 (8\%) and 0.11 (11\%) differ by 3 percentage points. 

\begin{figure}[!t]
    \centering
    \includegraphics[width=\linewidth]{Figures/CWE-epss.jpg}
    \caption{Distribution of absolute score differences (in percentage points) between EPSS and historical CWE-based exploitability scores across 2021--2025.}
    \label{fig:epss_cwe_difference}
\end{figure}
Figure~\ref{fig:epss_cwe_difference} shows the distribution of these differences across 2021--2025, with median differences ranging from approximately 1.1 to 1.4 percentage points, indicating relatively small numerical differences between the two approaches for the central portion of the evaluated CVEs.

Figure~\ref{fig:epss_cwe_difference} presents the distribution
of absolute score differences across the five evaluation years
(2021--2025). The median differences ranged from approximately
1.1 to 1.4 percentage points, with broadly comparable
interquartile ranges across the evaluation period.

To evaluate whether the numerical similarity between
historical CWE-based exploitability scores and EPSS
translates into comparable predictive performance,
we calculated PR-AUC and Precision@100 using the same
CVEs and exploitation labels for each evaluation year. Table~\ref{tab:cwe_epss_comparison} compares the predictive
performance of historical CWE-based exploitability scores
and EPSS across five evaluation years. EPSS consistently
outperformed the CWE-based approach, achieving a mean
PR-AUC of 0.3755 compared with 0.0436 and a mean
Precision@100 of 0.788 compared with 0.048.
Although historical CWE exploitation frequencies provide
a predictive signal, their vulnerability-level
discrimination is substantially weaker than EPSS in
this retrospective evaluation. Notably, EPSS scores
were obtained from an October 2026 snapshot, which
limits the temporal comparability of the two approaches.


\begin{table}[t]
\centering
\caption{Comparison of historical CWE-based exploitability
scores and EPSS for vulnerability prioritization (2021--2025).}
\label{tab:cwe_epss_comparison}
\small
\begin{tabular}{lcccc}
\hline
\textbf{Year} &
\multicolumn{2}{c}{\textbf{PR-AUC}} &
\multicolumn{2}{c}{\textbf{Precision@100}} \\
\cline{2-5}
& \textbf{CWE} & \textbf{EPSS} &
\textbf{CWE} & \textbf{EPSS} \\
\hline
2021 & 0.0533 & 0.4193 & 0.05 & 0.78 \\
2022 & 0.0472 & 0.3298 & 0.05 & 0.64 \\
2023 & 0.0408 & 0.3974 & 0.03 & 0.77 \\
2024 & 0.0410 & 0.3313 & 0.00 & 0.84 \\
2025 & 0.0355 & 0.3995 & 0.11 & 0.91 \\
\hline
\textbf{Mean} &
\textbf{0.0436} & \textbf{0.3755} &
\textbf{0.048} & \textbf{0.788} \\
\hline
\end{tabular}
\end{table}

\section{Literature Review}

Vulnerability prioritization has traditionally relied on
the Common Vulnerability Scoring System (CVSS), although
severity scores alone are insufficient indicators of
real-world exploitation. The Exploit Prediction Scoring
System (EPSS) developed by Jacobs et
al.~\cite{jacobs2023epss} addresses this limitation by
estimating the probability that an individual CVE will
be exploited. Subsequent studies have extended this
direction. Sherif et al.~\cite{sherif2026kri} combined
threat, impact, and exposure information into a Key
Risk Indicator framework. Khanmohammadi et
al.~\cite{khanmohammadi2025exploitabilitybirthmark}
introduced \textit{ExploitabilityBirthMark}, which
predicts vulnerability exploitability using information
available at disclosure. Jiang et
al.~\cite{jiang2025survey} surveyed the broader
development of dynamic vulnerability prioritization
approaches that incorporate exploitation-related
information. These studies primarily focus on
individual CVEs rather than systematically examining
whether the historical exploitation patterns of CWE
categories can provide an additional prioritization
signal.

A second research direction focuses on measuring and
ranking software weaknesses. Mell and
Gueye~\cite{mell2020suite} proposed metrics for
evaluating CWE categories based on vulnerability
frequency, severity, impact, and exploitability.
Galhardo et
al.~\cite{galhardo2020measurements}
investigated limitations of the MITRE CWE Top 25 \cite{cwe_top25} ranking methodology and proposed refinements to
better balance exploitability and impact. Although
these approaches help identify important weakness
categories, their primary objective is to characterize
and rank weaknesses in aggregate rather than estimate
the exploitation risk of newly disclosed CVEs.

A third research direction investigates automated
CVE-to-CWE mapping, which aims to identify the
underlying software weakness from a vulnerability
description. Existing approaches can be broadly
categorized into keyword-based methods, supervised
language models, generative LLMs, and semantic
embedding-based retrieval.

Keyword-based approaches identify weakness categories
using predefined terms and patterns, such as SQL
injection, cross-site scripting (XSS), and
cross-site request forgery (CSRF). While these
methods can be effective when vulnerability
descriptions explicitly mention recognizable
weaknesses, they are sensitive to terminology,
may produce false positives, and struggle with
descriptions that lack distinctive keywords.
Consequently, semantic language representations
have become increasingly important for automated
weakness classification.

Supervised transformer-based approaches formulate
CVE-to-CWE mapping as a classification task.
Das et al.~\cite{das2021v2w} proposed V2W-BERT,
a hierarchical BERT-based classification framework,
reporting up to 97\% accuracy under random data
partitioning and 94\% under temporal partitioning. Albanese et al. ~\cite{albanese2024cve2cwe} proposed CVE2CWE, which represents each CWE using
TF-IDF features extracted from historical CVE
descriptions and applies cosine similarity to
classify new vulnerabilities. They achieved
69.9\% top-1 accuracy across the 25 most frequent
CWEs, which decreased to 57.2\% when considering
50 CWEs. More recently, Mosievskiy et
al.~\cite{mosievskiy2026roberta} developed a
RoBERTa-based CWE classifier, reporting 87.4\%
top-1 accuracy on their benchmark. These findings
demonstrate the effectiveness of supervised
language representations for weakness
classification. However, conventional classifiers
typically operate over a predefined set of labels,
which can limit their ability to accommodate
previously unseen or infrequent CWE categories
without additional training.

Generative LLMs provide an alternative approach in
which vulnerability descriptions are supplied through
prompts and the model generates the corresponding CWE
identifier. Liu et al.~\cite{liu2023not}
investigated ChatGPT-based vulnerability mapping
using 13,513 CVEs, reporting exact-match accuracies
of 50.88\% and 53.13\% under two prompting
strategies. When hierarchical relationships were
considered, parent-level agreement increased to
72.20\% and 76.13\%, respectively. These results
highlight the difficulty of identifying precise
weakness identifiers, even when broader weakness
categories are correctly recognized.

Simonetto et al.~\cite{simonetto2026knowing}
further investigated CWE classification using
fine-tuned language models and hierarchical
information. Their study reported an approximately
five-percentage-point improvement in F1-score over
a supervised baseline. In their comparison,
the previously reported prompted ChatGPT approach
achieved approximately 61\% weighted F1-score,
compared with 72\% for a supervised
BiGRU-TextCNN baseline, while zero-shot GPT-3.5
achieved approximately 62\%. Their strongest
LLM-based approach required fine-tuning a
7-billion-parameter model, with reported training
costs of approximately 95 hours and 19.7 kWh,
compared with 0.9 kWh for BERT.

Despite their flexibility, generative approaches
introduce several challenges. First, large models
can incur substantial computational and inference
costs relative to smaller task-specific encoders.
Second, fine-tuned classification models may be
restricted to a fixed label space. For example,
Simonetto et al.~\cite{simonetto2026knowing}
focused on the 57 most frequent CWE categories,
which covered approximately 95\% of their labeled
CVEs, potentially limiting coverage of rare or
newly introduced weaknesses. Third, evaluating
generative LLMs presents a potential data
contamination challenge because publicly available
CVE descriptions and CWE assignments may have
appeared in their pretraining corpora. To mitigate
this concern, Simonetto et al. constructed a
separate evaluation set of 2,777 CVEs published
after the release of their base model. These
limitations motivate the investigation of
computationally efficient approaches that can
generalize to newly disclosed vulnerabilities.

Embedding-based methods offer another alternative
by representing CVE descriptions and CWE definitions
in a shared semantic space and ranking candidate
weaknesses according to their similarity.
Haddad et al.~\cite{haddad2023automated}
formulated CVE-to-CWE mapping as a semantic ranking
problem, evaluating Sentence-BERT and rankT5
models on a manually annotated dataset of 4,012
CVEs. Their results demonstrated improvements
over traditional information retrieval and
transformer-based baselines. Unlike fixed-label
classifiers, bi-encoder retrieval architectures
can evaluate candidate CWE descriptions without
requiring a dedicated output neuron for each
weakness. Moreover, similarity scores provide
a relative measure for ranking alternative CWE
assignments, although these scores should not
be interpreted as calibrated classification
probabilities.

Despite these advances, several challenges remain,
including distinguishing semantically overlapping
CWEs, handling imbalanced weakness distributions,
maintaining performance on newly disclosed
vulnerabilities, and accommodating changes to the
CWE taxonomy. Existing approaches also differ
substantially in their computational requirements,
dependence on labeled data, and ability to exploit
hierarchical relationships among weaknesses.

% To address these challenges, we propose HiCWE
% (Hierarchy-Aware Contrastive CWE Encoder), a
% lightweight bi-encoder framework based on
% \texttt{all-MiniLM-L6-v2} and fine-tuned using
% Low-Rank Adaptation (LoRA). HiCWE incorporates
% the CWE hierarchy and contrastive learning with
% hard negatives to improve semantic alignment
% between CVE descriptions and CWE definitions.
% Unlike generative LLMs, the framework retrieves
% predictions from an explicitly defined set of
% candidate weaknesses, preventing the generation
% of invalid CWE identifiers. Its retrieval-based
% architecture also enables newly introduced CWE
% definitions to be considered without modifying
% a fixed classification output layer. We evaluate
% the model using temporally separated CVE records
% to assess its generalization to future vulnerability
% disclosures.

Existing studies address CVE-level exploitability
prediction, CWE risk ranking, and automated
CVE-to-CWE mapping largely independently.
However, they have not systematically examined
whether historical CWE exploitation patterns
remain stable enough to support early vulnerability
prioritization. We address this gap by combining
historical CWE exploitability analysis with
hierarchy-aware CVE-to-CWE mapping in a unified
prioritization framework.

\section{Threats to Validity}

Like all empirical studies, this work has several limitations. Our analysis assumes that CWE assignments in the National Vulnerability Database (NVD) accurately represent the underlying software weaknesses. Although Finding 2 showed that few CWE mappings change after publication, some may still be incomplete, overly coarse, or incorrect. Similarly, exploitability was identified using CISA KEV, VulnCheck, and GreyNoise. While combining these sources improves coverage, exploited vulnerabilities absent from all three datasets may be incorrectly classified as non-exploited. Moreover, our LLM evaluation measures CWE classification accuracy rather than the end-to-end performance of a vulnerability prioritization system.

Methodological choices may also influence the results. Historical exploitability analyses considered only CWE-year combinations with sufficient observations, and exploitability rankings were restricted to CWEs satisfying minimum sample-size and statistical significance requirements. Although sensitivity analyses showed that the identified CWEs remain largely consistent under different thresholds, alternative thresholds or ranking weights may produce different rankings. Likewise, the LLM evaluation used specific prompting strategies, taxonomy configurations, and a limited set of open-source models; different prompts, models, or retrieval methods may yield different results.

Our empirical study relies on publicly disclosed CVEs from the NVD and publicly available exploitation datasets. Consequently, undisclosed vulnerabilities or exploits known only through proprietary threat intelligence are not represented, limiting generalizability to private vulnerability repositories or organization-specific environments. Similarly, the evaluated LLM framework targets software weaknesses defined by the CWE taxonomy and may require adaptation for other vulnerability taxonomies.

Finally, although we observed statistically significant temporal stability of historical exploitability at the CWE level, the correlations are moderate rather than strong. Historical exploitability should therefore be viewed as one informative signal rather than a complete predictor of exploitation. It is intended to complement other factors, including CVSS metrics, affected products, exploit availability, and threat intelligence, in future vulnerability prioritization frameworks.


\section{Conclusion and Future Work}
\label{sec:conclusion}

In this paper, we investigated the use of historical
CWE exploitability patterns and automated CWE classification
for early vulnerability prioritization. Our empirical
analysis of 192,788 CVEs revealed significant differences
in exploitation rates across CWE categories and moderate
temporal stability (mean Spearman's $\rho=0.480$),
demonstrating the potential of historical weakness-level
information as a prioritization signal.

We introduced HiCWE, a hierarchy-aware contrastive
bi-encoder for automated CVE-to-CWE mapping, achieving
75.3\% top-1 classification accuracy and outperforming
the evaluated baseline classifiers. We further demonstrated
that historical CWE exploitation rates can prioritize
future vulnerabilities, achieving ROC-AUC scores between
0.692 and 0.727 across temporal evaluations from
2021 to 2025 using official NVD CWE assignments.

Future work will focus on evaluating the complete
framework by incorporating HiCWE-predicted CWE labels
into the exploitability ranking process. We also plan
to investigate multi-label CWE classification, improve
generalization to rare and newly introduced weaknesses,
and integrate historical CWE exploitability with
complementary indicators such as CVSS, EPSS, and
threat intelligence to improve prioritization accuracy.

Overall, our findings demonstrate that historical
CWE exploitability provides a useful complementary
signal for early vulnerability prioritization and
that hierarchy-aware semantic learning can improve
automated weakness identification.



\clearpage
\bibliographystyle{splncs04}
\bibliography{references.bib}

\clearpage
\section*{Appendix}
\begin{figure}[H]
\centering
\includegraphics[width=1\columnwidth]{Figures/model/overall.png}
\caption{CVE$\rightarrow$CWE Accuracy@$k$ for the untuned baseline, the two finetuned recipes, and the published classifier. One panel per target.}
\label{fig:accuracy}
\end{figure}

\begin{table*}[ht]
\centering
\caption{Temperature sensitivity of the evaluated LLMs under Config C ($n=150$).}
\label{tab:temperature_all}
\small
\begin{tabular}{llcccc}
\toprule
\textbf{Model} & \textbf{Metric} &
\textbf{T=0.0} &
\textbf{T=0.3} &
\textbf{T=0.7} &
\textbf{T=1.0} \\
\midrule

\multirow{4}{*}{qwen2.5-coder (14B)}
& Exact-match accuracy  & 23.3\% & 26.0\% & 25.3\% & 19.3\% \\
& Parent-level accuracy & 28.0\% & 26.7\% & 30.7\% & 22.7\% \\
& Prediction diversity  & 38/150 & 45/150 & 46/150 & 57/150 \\
& Avg.\ time (s/CVE)    & 3.9 & 3.9 & 3.8 & 3.8 \\
\midrule

\multirow{4}{*}{qwen2.5 (14B)}
& Exact-match accuracy  & 17.3\% & 14.0\% & 18.0\% & 16.7\% \\
& Parent-level accuracy & 20.7\% & 16.7\% & 21.3\% & 20.0\% \\
& Prediction diversity  & 47/150 & 47/150 & 48/150 & 55/150 \\
& Avg.\ time (s/CVE)    & 3.8 & 3.7 & 3.6 & 3.6 \\
\midrule

\multirow{4}{*}{phi4 (14B)}
& Exact-match accuracy  & 16.7\% & 16.7\% & 16.7\% & 16.7\% \\
& Parent-level accuracy & 16.0\% & 18.0\% & 18.0\% & 16.7\% \\
& Prediction diversity  & 34/150 & 42/150 & 51/150 & 52/150 \\
& Avg.\ time (s/CVE)    & 5.6 & 5.3 & 5.3 & 5.4 \\
\midrule

\multirow{4}{*}{llama3.2 (3B)}
& Exact-match accuracy  & 7.3\% & 5.3\% & 3.3\% & 2.0\% \\
& Parent-level accuracy & 9.3\% & 6.7\% & 4.7\% & 2.0\% \\
& Prediction diversity  & 30/150 & 38/150 & 64/150 & 78/150 \\
& Avg.\ time (s/CVE)    & 1.2 & 1.2 & 1.2 & 1.1 \\
\bottomrule
\end{tabular}
\end{table*}



\end{document}
